\documentclass[11pt]{article}
% Shared preamble for the three living documents.
\usepackage[T1]{fontenc}
\usepackage{lmodern}
\usepackage[utf8]{inputenc}
\usepackage[margin=2.4cm]{geometry}
\usepackage{microtype}
\usepackage{amsmath,amssymb}
\usepackage{siunitx}
\usepackage{booktabs}
\usepackage{enumitem}
\usepackage{xcolor}
\usepackage{tikz}
\usetikzlibrary{arrows.meta,positioning,calc}
\usepackage{listings}
\usepackage[backend=biber,style=numeric,sorting=none]{biblatex}
\addbibresource{references.bib}
\usepackage[colorlinks=true,linkcolor=blue!50!black,citecolor=green!40!black,urlcolor=blue!60!black]{hyperref}

\newcommand{\vect}[1]{\mathbf{#1}}
\newcommand{\unitv}[1]{\hat{\mathbf{#1}}}
\newcommand{\dd}{\mathrm{d}}
\newcommand{\project}{Physics Sandbox}

\lstset{basicstyle=\ttfamily\small,columns=flexible,keepspaces=true,breaklines=true,frame=single,framerule=0.3pt}
\setlist{itemsep=2pt}
\setlength{\emergencystretch}{3em}

\title{\project{} --- Math Notes}
\author{Living document}
\date{Last updated: \today}

\begin{document}
\maketitle
\noindent This document describes the mathematics and numerical methods behind
the simulation: how continuous equations become the discrete updates the code
actually performs, and what errors that introduces. Each implementation step
adds a section.

\tableofcontents

\section{Notation}
Vectors are bold: $\vect{r} = (x,y,z)$, with magnitude $|\vect{r}| = \sqrt{x^2+y^2+z^2}$.
Subscript $h$ means the horizontal part, $\vect{v}_h = (v_x, 0, v_z)$. Discrete times are
$t_n = n\,\Delta t$, and $\vect{r}_n$ means $\vect{r}(t_n)$ as computed by the program.

\section{Milestone 1: a ball on a table}

\subsection{Direction from two angles}
A direction is the point on the unit sphere with azimuth $\theta$ and elevation
$\varphi$ (spherical coordinates with $y$ as the polar axis):
\begin{equation}
  \unitv{n} = (\cos\varphi\cos\theta,\ \sin\varphi,\ -\cos\varphi\sin\theta), \qquad |\unitv{n}| = 1.
\end{equation}
The minus sign on $z$ makes $\theta$ increase counter-clockwise when the $x$--$z$ plane is
seen from above ($+y$), because $+z$ points toward the viewer in a right-handed, $y$-up frame.

\subsection{The equation to solve}
Newton's law is a second-order ODE, written as a first-order system for the state $(\vect{r},\vect{v})$:
\begin{equation}
  \dot{\vect{r}} = \vect{v}, \qquad \dot{\vect{v}} = \vect{a}(\vect{r},\vect{v},t) = \frac{\vect{F}_\text{net}(\vect{r},\vect{v},t)}{m}.
\end{equation}
The right-hand side is only \emph{piecewise} smooth: it jumps when a force switches on or off,
when friction changes from static to kinetic, and when the ball leaves or lands on a surface.

\subsection{Time grid}
Time is always computed as $t_n = n\,\Delta t$ rather than by repeatedly adding $\Delta t$.
Repeated addition accumulates rounding error (e.g.\ $0.1$ has no exact binary representation),
so after many steps $\sum \Delta t \ne n\Delta t$~\cite{goldberg1991}. Multiplying keeps every
time reproducible, which the timeline relies on when seeking. The step size is limited to
$\SI{0.001}{s}\le\Delta t\le\SI{0.1}{s}$.

\subsection{Integrator: semi-implicit (symplectic) Euler}
The update used is~\cite{hairer2006}
\begin{align}
  \vect{v}_{n+1} &= \vect{v}_n + \vect{a}(\vect{r}_n,\vect{v}_n,t_n)\,\Delta t, \\
  \vect{r}_{n+1} &= \vect{r}_n + \vect{v}_{n+1}\,\Delta t.
\end{align}
It differs from \emph{explicit} Euler only in using the \emph{new} velocity to update the position.
Both are first-order accurate (global error $O(\Delta t)$)~\cite{press2007}, but semi-implicit Euler is
\emph{symplectic}: for conservative systems such as a pendulum or an orbit, its energy error
stays bounded instead of growing steadily~\cite{hairer2006}. That matters for later milestones.

\paragraph{Exact error for a constant force.} With constant $\vect{a}$ and $\vect{v}_0 = 0$:
$\vect{v}_n = n\,\vect{a}\,\Delta t$ and
\begin{equation}
  \vect{r}_n - \vect{r}_0 = \sum_{k=1}^{n} \vect{v}_k\,\Delta t = \vect{a}\,\Delta t^2\,\frac{n(n+1)}{2}
  = \tfrac12\vect{a}\,t_n^2 + \tfrac12\vect{a}\,\Delta t\, t_n .
\end{equation}
The velocity is exact, and the position overshoots the exact $\tfrac12 a t^2$ by
$\tfrac12 a\,\Delta t\,t$. This is first-order in $\Delta t$, and the automated test checks it.
For $a = \SI{2}{m/s^2}$, $\Delta t = \SI{0.01}{s}$ and $t = \SI{2}{s}$, the error is \SI{2}{cm}
out of \SI{4}{m}, i.e.\ 0.5\,\%.

\subsection{Friction is non-smooth}
The kinetic friction term $-\mu_k N\,\vect{v}_h/|\vect{v}_h|$ is discontinuous at $\vect{v}_h=\vect{0}$, and
static friction is \emph{set-valued} there: it can be anything with magnitude up to $\mu_s N$~\cite{stewart2000}.
A naive integrator makes a stopping ball jitter back and forth. Two rules avoid that:
\begin{enumerate}
  \item \textbf{Rest test.} If $|\vect{v}_h| < \varepsilon_v = \SI{e-6}{m/s}$, the ball counts as at rest and the
        static rule applies: friction exactly cancels the horizontal push if $|\vect{F}_h|\le\mu_s N$.
  \item \textbf{No reversal.} If, while sliding, a step would make $\vect{v}_{h,n+1}\cdot\vect{v}_{h,n} \le 0$, then
        friction would have stopped the ball inside the step, so $\vect{v}_{h,n+1}$ is set to $\vect{0}$.
\end{enumerate}
The stopping time is then correct to within one step, which is also tested.

\subsection{Contact and landing}
The ball is \emph{in contact} with a surface of height $y_s$ if
$|y - (y_s + r)| \le \varepsilon_c = \SI{e-6}{m}$, $v_y \le \varepsilon_v$, and the centre lies within the
surface's bounds. While in contact, $N$ cancels the downward external force, so $a_y = 0$.

\emph{Landing} is detected when a step carries the centre from above the contact height to
below it: $y_n \ge y_s + r - \varepsilon_c$ and $y_{n+1} < y_s + r$. The position is snapped to
$y_s + r$ and the downward velocity is removed. The landing time therefore has an error of at most
$\Delta t$. Detecting the exact crossing time (event location) is a possible later refinement~\cite{press2007}.

\subsection{Recorded quantities}
At each step the program stores $t_n$, $\vect{r}_n$, $\vect{v}_n$ and the acceleration $\vect{a}_n = \vect{F}_\text{net}(\vect{r}_n,\vect{v}_n,t_n)/m$
\emph{at that state}. It also stores these derived values:
\begin{align}
  \Delta\vect{r}_n &= \vect{r}_n - \vect{r}_0, \\
  s_{n+1} &= s_n + |\vect{v}_{n+1}|\,\Delta t, \qquad s_{x,n+1} = s_{x,n} + |v_{x,n+1}|\,\Delta t \ \ (\text{same for } y, z).
\end{align}
Because $\vect{r}_{n+1}-\vect{r}_n = \vect{v}_{n+1}\Delta t$ exactly, $s_n$ is the \emph{exact} length of the
polyline through the computed positions. It is a right-endpoint Riemann sum for
$\int |\vect{v}|\,\dd t$, consistent with the integrator. By the triangle inequality,
$s_n \ge |\Delta\vect{r}_n|$.

\subsection{Determinism}
Given the same parameters, forces and $\Delta t$, the computation is a pure function of its inputs, so
the results are bit-for-bit identical. Rewinding just reads stored samples. Changing a past input
recomputes from $t=0$; applying a new force at time $t_k$ keeps samples $0..k$ and recomputes the rest.

\section{Milestone 1.1: interaction geometry}

\subsection{From a vector back to magnitude and angles}
Dragging the arrow produces a force vector $\vect{F}$. The panel shows it in the same
spherical coordinates as before, by inverting the direction formula:
\begin{equation}
  F = |\vect{F}|,\qquad
  \varphi = \arcsin\!\left(\frac{F_y}{F}\right),\qquad
  \theta = \operatorname{atan2}(-F_z,\ F_x)\ \ (\text{mapped to } [0^\circ, 360^\circ)).
\end{equation}
$\operatorname{atan2}$ resolves the quadrant that $\arctan(-F_z/F_x)$ cannot. When the force is
vertical ($F_x=F_z=0$) the azimuth is undefined, so the previous value is kept. Values are rounded
to \SI{0.01}{N} and \ang{0.1}.

\subsection{Turning a 2D drag into a 3D vector}
A mouse or finger position fixes only a \emph{ray} from the camera into the scene. The ray is
turned into a point by intersecting it with a plane:
\begin{itemize}
  \item by default, the plane through the arrow tip that faces the camera (normal $= $ the camera's
        viewing direction). In the top view this plane is horizontal; in the front and side views it is vertical;
  \item with Shift held, the horizontal plane through the tip.
\end{itemize}
For a ray $\vect{p}(s) = \vect{o} + s\,\unitv{d}$ and a plane $\unitv{n}\cdot\vect{x} = c$, the
intersection is at $s = (c - \unitv{n}\cdot\vect{o})/(\unitv{n}\cdot\unitv{d})$ (no intersection if
$\unitv{n}\cdot\unitv{d}=0$). The force is then $\vect{F} = (\vect{x}_\text{tip} - \vect{r}_\text{ball})/k$
with $k = \SI{0.06}{m/N}$.

\subsection{Axis gizmo}
The corner gizmo shows where the world axes point as seen by the camera. If $q$ is the camera's
orientation (a unit quaternion), a world direction $\unitv{e}$ has camera-space coordinates
$\vect{v} = q^{-1}\,\unitv{e}\,q$. The gizmo draws $\unitv{e}$ at screen offset $(v_x, -v_y)$
times a fixed radius, and paints the axes in order of increasing $v_z$ (farthest first), so the
axis pointing toward the viewer is drawn on top.

\section{Milestone 1.2: picking arrows and drawing an endless grid}

\subsection{Which arrow is under the pointer?}
Each arrow is projected to the screen as a segment from $\vect{a}$ (tail) to $\vect{b}$ (tip). The
distance from the pointer $\vect{p}$ to the segment is
\begin{equation}
  d = \left|\vect{p} - (\vect{a} + t^\ast(\vect{b}-\vect{a}))\right|,\qquad
  t^\ast = \operatorname{clamp}\!\left(\frac{(\vect{p}-\vect{a})\cdot(\vect{b}-\vect{a})}{|\vect{b}-\vect{a}|^2},\,0,\,1\right),
\end{equation}
i.e.\ the distance to the closest point of the segment. The arrow with the smallest $d$ under a
threshold (\SI{8}{px} for a mouse, \SI{22}{px} for touch) is chosen. The tooltip's angles use the same
inversion as in Milestone~1.1: $\theta=\operatorname{atan2}(-F_z, F_x)$, $\varphi=\arcsin(F_y/|\vect{F}|)$.

\subsection{An anti-aliased procedural grid}
The grid is computed for each pixel from that pixel's world coordinates $(x,z)$ rather than drawn as
geometry. For spacing $s$, let $\vect{c} = (x,z)/s$. The distance to the nearest grid line, in
units of whole grid cells, is $|\operatorname{fract}(c-\tfrac12)-\tfrac12|$. Dividing by
$\operatorname{fwidth}(c) = |\partial c/\partial X| + |\partial c/\partial Y|$, the change of $c$
from one screen pixel to the next, converts it to \emph{pixels}. That gives lines of constant screen
width that stay anti-aliased at any distance.

Far away, many grid lines fall within one pixel and would alias into moiré patterns. The grid's
opacity is therefore faded to zero with distance $D$ from the camera using a smooth step,
\[
  \operatorname{fade}(D) = 1 - \operatorname{smoothstep}(D_0/4,\ D_0,\ D),\qquad D_0 = \max(40,\ 12\,|y_\text{cam}|)\ \text{m},
\]
and linear fog blends the floor itself into the background colour between \SI{60}{m} and \SI{400}{m}. Together
these make the finite floor plane (\SI{10}{km} across) look endless.

\paragraph{Colour space.} Lit materials compute in \emph{linear} RGB, and the result is converted to
sRGB for display. Grid colours written in sRGB must therefore be converted first
($c_\text{lin} \approx c_\text{sRGB}^{2.2}$). Otherwise mid-grey lines come out almost the same
colour as the floor.

\section*{Changelog}
\begin{itemize}
  \item \textbf{Milestone 1} --- semi-implicit Euler, friction regularisation, landing snap,
        path-length accumulation, exact time grid.
  \item \textbf{Milestone 1.1} --- inverting the direction formula, ray--plane picking for
        arrow dragging, the axis gizmo projection.
  \item \textbf{Milestone 1.2} --- point-to-segment picking, derivative-based grid lines, distance fade.
\end{itemize}

\printbibliography
\end{document}
